% !PRINTED 582
\weberpage{582}
Ueber die Frage, welche der drei Wurzeln der cubischen
Gleichung (19) für $\xi = 3\eta + 1$, welche für $3\eta_1 + 1$ oder für
$3\eta_2 + 1$ zu setzen ist, lässt sich allgemein so viel sagen, dass
das Vorzeichen des Productes (17) mit dem Vorzeichen von (18),
das durch die Definition (1), (2) bestimmt ist, freilich aber noch
von der Wahl der primitiven Wurzel $g$ abhängt, übereinstimmen
muss. Dadurch sind von den sechs möglichen Zuordnungen drei
ausgeschlossen. Welche Zuordnung unter den drei übrigen zu
treffen ist, hängt von der Wahl von $r$ ab, und würde zu ähnlichen
Untersuchungen Anlass geben, wie sie im vorigen Paragraphen zur
Bestimmung des Zeichens der Quadratwurzel durchgeführt sind\footnote{\gesperrt{Kummer}, Journ. f. Mathematik, Bd. 32.}.

Wir heben noch den durch die Formeln (14), (16) bewiesenen
Satz hervor:

\gesperrt{Ist} $n$ \gesperrt{eine Primzahl von der Form} $3k + 1$, \gesperrt{so ist} $n$
\gesperrt{durch die Form} $a^2 - ab + b^2$ \gesperrt{und} $4n$ \gesperrt{durch die Form}
$A^2 + 27B^2$ \gesperrt{darstellbar, wo} $a, b, A, B$ \gesperrt{ganze Zahlen sind.}

Daraus ergiebt sich noch, dass $n$ auch in der Form $x^2 + 3y^2$
darstellbar ist. Denn wenn von den beiden Zahlen $a, b$ eine,
etwa $a$ gerade ist, so ist
\[
n = \left(\frac{a}{2} - b\right)^2 + \tfrac{3}{4}a^2,
\]
und wenn $a$ und $b$ beide ungerade sind, so ist
\[
n = \tfrac{1}{4}(a + b)^2 + \tfrac{3}{4}(a - b)^3. % [sic] printed exponent 3, expected 2
\]

Auch die zweigliedrigen Perioden der $9^{\mathrm{ten}}$ Einheitswurzeln
genügen einer cubischen Gleichung, und weil diese Gleichung
bei der später zu behandelnden allgemeinen Theorie der cubischen
Kreistheilungskörper eine wichtige Rolle spielt, wollen wir
sie der Vollständigkeit halber hier betrachten, obwohl sie eigent\-
lich in ein allgemeineres Gebiet gehört, in dem die Grade der
Einheitswurzeln keine Primzahlen mehr sind.

Eine $9^{\mathrm{te}}$ Einheitswuzel $r$ genügt der Gleichung $6^{\mathrm{ten}}$ Grades % [sic] \textquotedbl{}Einheitswuzel\textquotedbl{} as printed
\weq{21}{r^6 + r^3 + 1 = 0,}
und wir haben drei conjugirte zweigliedrige Perioden
\weq{22}{\eta = r + r^{-1}, \quad \eta' = r^2 + r^{-2}, \quad \eta'' = r^4 + r^{-4},}
die die Wurzeln der cubischen Gleichung
\weq{23}{\eta^3 - 3\eta + 1 = 0}
%CONT>>
